\section{Relation to Other Forms of Slicing \& Type Debugging}
\label{sec:comparisons}
\label{sec:related-work}

Program slicing methods have been explored extensively, in various forms and for many applications~\cite{TipSurvey1995}. Here we compare type slicing with these, being loosely related in foundation to \emph{Galois slicing}~\cite{Perera2012}, and in application to \emph{type error slicing}~\cite{Wand1986,HaackWells2003}, before relating to the broader type debugging literature (\S\ref{sec:cmp-debuggers}).

\subsection{Classical Program Slicing}
\label{sec:cmp-program-slicing}

    Classical program slicing~\cite{TipSurvey1995} was first introduced by Weiser~\cite{Weiser1984}. Here, a slice is a subprogram preserving the value of a chosen variable at a chosen program point: a `criterion'. Similarly to type slicing, the subprogram is a syntactically well-formed program which can be analysed (evaluated) accordingly.

The criterion is superficially similar to type slicing queries, with two significant differences. First, a type query is itself a slice of a type, which has \textit{compound} structure, rather than a single variable, and type slices cannot in general \textit{preserve} all queries exactly (\cref{lem:exact-not-exist}), so we ask only for \textit{validity} (\cref{def:syn-slice}), which remains a `sufficient' explanation. 

Classical slicing distinguishes \textit{static} slices (preserved over all executions) from \textit{dynamic} slices (over one fixed execution)~\cite{KorelLaski1988}. Type slicing can be viewed as a \textit{dynamic} slicing method, as it fixes a particular typing derivation; however, most type systems, including ours, have no non-trivial non-determinism, making this distinction largely irrelevant.

\subsection{Galois Slicing}
\label{sec:cmp-galois}

Galois slicing~\cite{Perera2012} is a lattice-theoretic approach to dynamic program slicing that has been applied and related to a wide range of calculi and user interfaces~\cite{Ricciotti2017,PereraGarg2016,PereraVis2022,AtkeyPerera2025}. Type slicing shares the foundational use of lattices and precision, but differs decisively in the lack of a \textit{Galois connection}. 

In Perera's setting, a program evaluation lifts to a \textit{meet-preserving} map $\mathit{eval}_{\rho,e} : \lfloor \rho, e\rfloor \to \lfloor v\rfloor$ between the slice lattices of partial programs and partial values. By an adjoint functor theorem~\cite[Prop.~7.34]{DaveyPriestley}, such a map is the upper adjoint of a unique Galois connection, whose lower adjoint $\mathit{uneval}$ is the `backward slice function' sending a partial-value query to a \textit{least} program slice evaluating to it.

For total type systems assigning unique types like the core calculus (\S\ref{sec:core-calculus}), the natural type-slicing analogue of $\mathit{eval}_{\rho,e}$ maps partial assumptions and expressions to partial types by type checking.
\[
  \phi_D : \lfloor \Gamma, e \rfloor \to \lfloor \tau \rfloor
\]

However, this `forward map' $\phi_D$ is \textit{not meet-preserving} in general. Consider $D : \synthesis{(\mathbf{case}\;e\;\mathbf{of}\;x_1.\,e_1 \mid x_2.\,e_2)}{\tau}$ in which both branches independently synthesise the same type $\tau \sqsupset \gap$, and take the two slices $s_1$ keeping only the first branch and $s_2$ keeping only the second. Both are defined, yielding $\phi_D(s_1) = \phi_D(s_2) = \tau$, but their meet $s_1 \sqcap s_2$ omits both branches, giving $\phi_D(s_1 \sqcap s_2) = \gap \neq \tau = \phi_D(s_1) \sqcap \phi_D(s_2)$.

\subsection{Type Error Slicing}
\label{sec:cmp-tes}

Type error slicing, like type slicing, is a static analysis on a typing problem. Both share the broad goal of attributing type-level information back to source positions, but differ in language target, the underlying mathematical structure, and the scope.

Wand~\cite{Wand1986} and Haack and Wells~\cite{HaackWells2003} work in a globally inferred Hindley-Milner system; the program is reduced to a system of equality constraints over type variables, and a type error slice is computed as a minimal unsatisfiable subset of the constraint set, projected back to the corresponding minimal subset of source positions. Skalpel~\cite{Skalpel2017} matures and scales this approach to full Standard ML.

These slices are not structure-preserving programs but flat sets of constraints and sources, so the results cannot themselves be type checked or executed.

The scope also differs: type error slicing answers \emph{``why does this program fail to type-check?''}, whereas type slicing answers \emph{``why does this expression have this type?''}. Via marking (\S\ref{sec:marking})~\cite{Marking2024}, which elaborates every well-formed expression into a totally typed derivation, type slicing applies also to ill-typed programs. So, here, the question of ``why this type?'' subsumes ``why this fails to type-check?''.

\subsection{Type-Directed Slicing}
\label{sec:cmp-type-directed}

\citet{TypeDirectedSlicing} recently independently introduced \emph{type-directed slicing}, implemented for Java in the tool \textsc{Specimin}. Their motivation and application is more specific than ours: their slices serve the \emph{maintainers of a typechecker}, reducing a large program exhibiting a known bug to a small test case reproducing it.

Given a target program element and a typechecker specified by its type rules, a type-directed slicer statically computes a smaller program on which the typechecker behaves identically at the target: essentially type slicing restricted to maximal queries (no partial queries). However, the approach, theory, and proofs differ greatly, as do the languages considered, with ours extending to bidirectionally inferred types, applying to programs with fewer annotations and better support for first-class functions. On the other hand, their work applies in practice to a wider range of analyses, including exceptions, crashes, and nullness analysis.

Their work was empirically evaluated on historical bugs of \texttt{javac}, NullAway~\cite{NullAway}, and the Checker Framework~\cite{CheckerFramework}, suggesting our theory is similarly applicable to such real world use cases; conversely, partial queries may be a worthwhile extension to their tools, especially when a type checker produces multiple cascading bugs.

\subsection{Type Debuggers \& Type Explanation}
\label{sec:cmp-debuggers}

Beyond slicing, there is a broad literature on explaining, localising, and repairing type errors. The oldest line \emph{explains} types: Beaven and Stansifer~\cite{BeavenStansifer1993} and Duggan and Bent~\cite{DugganBent1996} generate textual explanations from instrumented unification, \citet{YangMichaelsonTrinder2002} explain inferred polymorphic types of well-typed programs, and \citet{Chitil2001} navigates compositional explanation graphs via algorithmic debugging. These systems answer the same question as a synthesis slice---why does this term have this type?---but produce prose or inference traces that grow with the program, whereas a type slice is itself a minimal well-formed program that can be re-checked and further refined. Interactive debuggers similarly answer such queries: Chameleon~\cite{StuckeyChameleon2003} over a constraint store for Haskell, and Argus~\cite{Argus2025} for Rust trait errors in the IDE.

A larger body of work improves error \emph{reports} for globally inferred languages: constraint-solving heuristics in Helium~\cite{HeerenHageSwierstra2003}, ranking likely error causes over constraint graphs~\cite{ZhangMyers2014,SHErrLoc2015}, minimum error sources via MaxSMT~\cite{PavlinovicKingWies2014,PavlinovicICFP2015}, black-box correcting sets~\cite{Mycroft2016}, subtyping-constraint flow messages~\cite{FlowMessages2023}, dynamic witnesses~\cite{SeidelJhalaWeimer2016}, and learned localisation~\cite{NATE2017} and repair~\cite{Rite2020}. These techniques answer ``where is the error?'' or ``how might it be fixed?'' rather than ``why this type?''; empirically, the true fix often lies outside the reported location~\cite{StudentTypeErrorFixes}, which type slicing surfaces directly via analysis and full slices (\S\ref{sec:full-slices}). As slices are themselves programs, these tools compose with type slicing (\S\ref{sec:intro}).


%%% Local Variables:
%%% mode: LaTeX
%%% TeX-master: "PolymorphicTypeSlicing"
%%% End:
